\chapter{Datasets and Preprocessing}
\thispagestyle{empty}

Two public datasets are used, one per disorder. They are referred to as DS-1 (IDD) and DS-2
(ADHD). Table~\ref{tab:datasets} summarises them.

\begin{table}[H]
\caption{The two datasets. Window counts are for 5~s windows with 50\% overlap.}
\label{tab:datasets}\centering\small
\begin{tabular}{lp{5.3cm}p{5.3cm}}
\hline
 & DS-1 (IDD)~\cite{sareen2020data} & DS-2 (ADHD)~\cite{motie2020} \\
\hline
Participants & 7 with IDD, 7 typically developing controls (TDC) & 61 with ADHD, 60 normal controls (NC) \\
Age & IDD group \R{DSoneIddAgeMin}--\R{DSoneIddAgeMax} years; controls not given in the dataset files & 7--12 years \\
Device & EMOTIV EPOC+, 14 channels & 19 channels, 10--20 system \\
Sampling rate & 128~Hz & 128~Hz \\
Recording & 2~min rest and 2~min listening to music & visual attention task, length depends on the child \\
Windows & \R{DSoneRestNWin} per condition (\R{DSoneRestWinPerSubjMin} per recording) & \R{DStwoNWin} (\R{DStwoNWinDis} ADHD, \R{DStwoNWinCtrl} NC) \\
Windows per subject & \R{DSoneRestWinPerSubjMin} & \R{DStwoWinPerSubjMin}--\R{DStwoWinPerSubjMax} \\
\hline
\end{tabular}
\end{table}

\section{DS-1: intellectual and developmental disorder}
DS-1~\cite{sareen2020data} contains recordings of 14 participants, 7 with IDD and 7 typically
developing controls, made with a 14-channel EMOTIV EPOC+ headset at 128~Hz. The electrodes are
AF3, F7, F3, FC5, T7, P7, O1, O2, P8, T8, FC6, F4, F8 and AF4. Each participant was recorded
for two minutes at rest and for two minutes while listening to music. The dataset provides both
the raw recordings and cleaned recordings; the cleaned recordings were produced by the dataset
authors in EEGLAB with a 1--30~Hz band-pass filter, removal of line noise, independent component
analysis and automatic rejection of artefact components.

The two recording conditions give three classification settings, all with the same
participants: rest only, music only, and rest and music together. Each recording gives
\R{DSoneRestWinPerSubjMin} windows, so the rest and music settings have \R{DSoneRestNWin}
windows each (\R{DSoneRestNWinDis} IDD and \R{DSoneRestNWinCtrl} TDC) and the combined setting
has \R{DSoneRestMusicNWin}.

\textbf{The participants of DS-1 are adults.} The metadata file distributed with the dataset
gives the ages of the seven IDD participants as \R{DSoneIddAgeMin}--\R{DSoneIddAgeMax} years; it
gives no ages for the controls. The IDD results of this project therefore concern adults, and
only the ADHD results concern children. In this report the terms ``child-level'' and
``subject-level'' are used interchangeably for quantities computed per person; for DS-1 the
person is an adult participant.

\section{DS-2: attention-deficit/hyperactivity disorder}
DS-2~\cite{motie2020} contains recordings of 121 children aged 7--12 years, 61 with ADHD and 60
controls, boys and girls. According to the dataset description reproduced
in~\cite{chandela2025}, ADHD was diagnosed by an experienced psychiatrist using DSM-IV criteria,
and the control children had no history of psychiatric disorders, epilepsy or high-risk
behaviour. EEG was recorded from 19 electrodes of the 10--20 system (Fp1, Fp2, F3, F4, C3, C4,
P3, P4, O1, O2, F7, F8, T3, T4, T5, T6, Fz, Cz, Pz) at 128~Hz. T3, T4, T5 and T6 are the older
names of the positions now called T7, T8, P7 and P8.

The recording protocol was a visual attention task: the child was shown pictures of cartoon
characters and asked to count them, with between 5 and 16 characters per picture. Each picture
was shown until the child answered, so the length of a recording depends on the child's
performance. DS-2 gives \R{DStwoNWin} windows, \R{DStwoNWinDis} from ADHD children and
\R{DStwoNWinCtrl} from controls, with between \R{DStwoWinPerSubjMin} and \R{DStwoWinPerSubjMax}
windows per child.

\textbf{Medication.} The children with ADHD had been taking methylphenidate (Ritalin) for up to
six months; the control children were not medicated. Since stimulant medication changes the
EEG~\cite{clarke2002}, \emph{every ADHD result in this report is a comparison of medicated
children with ADHD against unmedicated controls}, and a classifier may detect the effect of the
drug rather than the disorder. This caveat is repeated wherever ADHD results appear.

\section{Preprocessing}
\textbf{DS-1.} The cleaned recordings distributed with the dataset are used as they are. The
original cleaning used MATLAB toolboxes that have no exact Python equivalent, so using the
distributed cleaned data is the closest match to earlier work on this dataset.

\textbf{DS-2.} The recordings are filtered with a fourth-order Butterworth band-pass filter
(0.5--60~Hz) and a 50~Hz notch filter, both applied forwards and backwards so that no phase
shift is introduced. They are then processed with the discrete wavelet transform, using the
Coiflet-3 wavelet at decomposition level 6, in two steps. The first removes the slowest
component, where ocular artefacts concentrate. The second shrinks the detail coefficients with
a threshold chosen by Stein's unbiased risk estimate (SURE)~\cite{donoho1995} and reconstructs
the signal. At 128~Hz the six detail levels cover 32--64, 16--32, 8--16, 4--8, 2--4 and 1--2~Hz.
The filter order and the details of the wavelet steps are implementation choices
(Appendix~A).

\section{Segmentation}
Every recording is divided into windows of 5~s (640 samples) with 50\% overlap, so consecutive
windows share half of their samples. Each window is labelled with the diagnosis of the person
it comes from. Windows are the unit of feature extraction and classification throughout; a
decision for a person is obtained by combining the decisions for that person's windows
(Section~\ref{sec:protocol}).

Two consequences of this segmentation are central to the rest of the report. Windows of the
same person are not independent examples, both because they overlap and because they share
everything that is specific to that person and that recording. And the number of independent
units for statistics is the number of people (14 and 121), not the number of windows.

\section{Properties of the data relevant to classification}
\label{sec:datachecks}
Before any classifier is discussed, three properties of the raw data are recorded here. They
were measured with a descriptive script (\texttt{experiments/data\_checks.py}) that trains no
model; their consequences are discussed in Chapter~7.

\textbf{Group difference in beta power (DS-1).} For each rest recording, the relative power in
the beta band (13--30~Hz, as a fraction of 1--45~Hz power, averaged over channels) lies in
\R{DSoneRestRelBetaIdd} for the IDD participants and in \R{DSoneRestRelBetaTdc} for the
controls. The two ranges do not overlap, so this single number separates the 14 rest recordings
completely (area under the ROC curve \R{DSoneRestRelBetaAUC}). For the music recordings the
ranges overlap (\R{DSoneMusicRelBetaIdd} against \R{DSoneMusicRelBetaTdc}, AUC
\R{DSoneMusicRelBetaAUC}).

\textbf{An outlying recording (DS-1).} The median standard deviation of the
\R{DSoneNRecordings} cleaned recordings is \R{DSoneMedianSD}~$\mu$V. One control rest recording
(\R{DSoneOutlierFile}) has \R{DSoneOutlierSD}~$\mu$V, \R{DSoneOutlierTimes} times the median,
with peaks of \R{DSoneOutlierMax}~$\mu$V and nearly the same value on all 14 channels, which
indicates a large artefact common to all electrodes that remains in the distributed cleaned
data. No two recordings are duplicates of each other (\R{DSoneIdenticalPairs} identical pairs;
the largest correlation between any two recordings is \R{DSoneMaxPairCorr}).

\textbf{Recording length (DS-2).} Because the task continued until the child answered, the
groups differ in recording length: \R{DStwoLenAdhd} windows on average for children with ADHD
against \R{DStwoLenCtrl} for controls. Recording length alone predicts the group with an AUC of
\R{DStwoLenAUC}.
